\documentclass[11pt]{letter}
\usepackage[T1]{fontenc}
\usepackage{lmodern}
\usepackage[margin=1in]{geometry}
\usepackage[hidelinks]{hyperref}
\IfFileExists{xurl.sty}{\usepackage{xurl}}{}
\urlstyle{same}

\signature{Yicheng Yang\\University of Illinois Urbana-Champaign\\yy85@illinois.edu}
\address{Yicheng Yang\\University of Illinois Urbana-Champaign\\601 E John Street\\Champaign, IL 61820\\USA}

\begin{document}
\begin{letter}{The Editors\\Finance Research Letters}

\opening{Dear Editors,}

I am submitting ``Do Prediction Market Prices Equal Probabilities? Evidence from a Play-Money Contrast'' for consideration at \textit{Finance Research Letters} as a stand-alone empirical note.

This note tests whether the price--probability wedge in prediction markets varies with the presence of financial stakes, exploiting a natural contrast between traded real-money platforms (Polymarket, Kalshi) and the play-money platform Manifold. Using a one-parameter Wang specification purely as a measurement device, I estimate platform-level wedges on 290{,}730 resolved contracts. Real-money venues yield positive wedges ($\hat\lambda \approx 0.16$--$0.19$), while Manifold yields a negative wedge ($\hat\lambda = -0.218$), with pairwise differences above 0.38 in magnitude and implied $z$-statistics above 11. The pattern is robust to alternative price measures, complement-invariant specification, and sample restrictions.

The note builds on the journal's existing tradition of applying the Wang transform to incomplete-market settings (e.g., Badescu, Cui, and Ortega, 2016) by extending the framework to prediction markets---a setting where the Wang transform's incomplete-markets interpretation is particularly natural, since binary event payoffs are generally not spanned by traded assets. To the best of my knowledge, this is the first application of the Wang transform to prediction markets.

The manuscript is concise (well under the 2{,}500-word limit), focuses on a single clean empirical contrast, and explicitly acknowledges the identification limits of the cross-platform design. A separate working paper (Yang, 2026; SSRN 6468338) develops the structural selection framework in greater depth; the current submission is a focused empirical letter that scopes strictly to the real-money versus play-money contrast, appropriate to the letter format.

Thank you for considering this submission.

\closing{Sincerely,}

\end{letter}
\end{document}
